\chapter{A Guide to Menopause}

\section*{Definition and Overview}
Menopause is the permanent cessation of menstruation, confirmed retrospectively after 12 consecutive months without a period. It represents the natural end of a woman's reproductive years, when the ovaries stop releasing eggs and markedly reduce their production of oestrogen and progesterone. The years leading up to the final period, during which hormone levels fluctuate and symptoms often begin, are known as perimenopause or the menopausal transition. Menopause can also be induced by surgery, such as removal of both ovaries, or by medical treatments that damage ovarian function. This chapter provides an overview of the physiological changes, common symptoms, and the treatment and self-care options available to manage this stage of life.

\section*{Epidemiology}
Menopause is a universal biological event that affects every woman who lives to midlife. The average age of natural menopause is around 51 years in Australia and comparable populations, with most women experiencing the transition between their mid-40s and mid-50s. Approximately 1 in 10 women experience premature menopause before the age of 40, or early menopause between 40 and 45 years, which may be natural or caused by surgery, chemotherapy, or other factors. Up to three-quarters of women report vasomotor symptoms such as hot flushes and night sweats during the transition, and for roughly one in four these symptoms are severe enough to interfere with daily functioning. With increasing female life expectancy, many women now spend a substantial proportion of their lives in the post-menopausal state, making the long-term health effects of oestrogen deficiency an important public health concern.

\section*{Aetiology and Pathophysiology}
The central event in menopause is the age-related depletion of ovarian follicles, the structures that produce eggs and the steroid hormones oestrogen and progesterone. As follicular numbers decline, ovulation becomes erratic and eventually ceases, and circulating oestrogen levels fall to a low, stable plateau.

\begin{itemize}
    \item \textbf{Natural ageing:} the ovaries' finite supply of eggs is gradually exhausted, reducing hormone output and leading to irregular cycles that ultimately stop
    \item \textbf{Surgical menopause:} bilateral oophorectomy causes an abrupt and often more severe onset of symptoms because hormone levels fall suddenly
    \item \textbf{Medical treatments:} chemotherapy, pelvic radiotherapy, and some endocrine therapies can damage or suppress ovarian function, sometimes permanently
    \item \textbf{Premature ovarian insufficiency:} ovarian failure before age 40, which may be autoimmune, genetic (for example in Turner syndrome), or of unknown cause
    \item \textbf{Effects of oestrogen deficiency:} low oestrogen affects the hypothalamus, destabilising temperature regulation and causing vasomotor flushes; it also thins the vaginal epithelium, reduces bone density, and influences mood and sleep pathways
\end{itemize}

\section*{Clinical Features}
The symptoms of menopause vary widely between women in type and severity. The hallmark symptoms are vasomotor: hot flushes and night sweats, which can disrupt sleep and lead to significant daytime fatigue.

\begin{itemize}
    \item Hot flushes and night sweats, the most common and characteristic symptoms
    \item Sleep disturbance, insomnia, and resulting tiredness and irritability
    \item Vaginal dryness, itching, and discomfort, often with dyspareunia and reduced libido
    \item Mood changes, anxiety, irritability, and low mood
    \item Joint and muscle aches, headaches, and a sensation of "brain fog" or poor concentration
    \item Menstrual irregularity, with periods becoming longer, shorter, lighter, or heavier before stopping
    \item Urinary symptoms such as urgency, frequency, and recurrent urinary tract infections
    \item Skin and hair changes, including dryness and thinning
\end{itemize}

\section*{Differential Diagnosis}
Although menopause is a normal physiological event, its symptoms overlap with several medical conditions that should be considered, particularly when they are atypical or severe.

\begin{itemize}
    \item \textbf{Thyroid disease:} hypothyroidism and hyperthyroidism can both cause fatigue, mood disturbance, weight change, and menstrual irregularity
    \item \textbf{Anaemia:} iron deficiency may produce tiredness, breathlessness, and palpitations
    \item \textbf{Depression or anxiety:} primary mood disorders can occur independently of the menopause
    \item \textbf{Chronic fatigue syndrome} and other causes of persistent exhaustion
    \item \textbf{Pregnancy:} possible in women who are still ovulating intermittently
    \item \textbf{Structural gynaecological causes of bleeding:} uterine fibroids, polyps, or rarely endometrial cancer, especially with heavy or post-menopausal bleeding
\end{itemize}

\section*{When to Seek Medical Care}
Women should consult a doctor if menopausal symptoms are troublesome or affecting quality of life, if periods become very heavy or unpredictable, if bleeding occurs more than 12 months after the last period, or if they wish to discuss the risks and benefits of hormone therapy.

\textbf{Call 000 (or local emergency services) for:}
\begin{itemize}
    \item Very heavy vaginal bleeding that soaks pads or tampons within an hour, especially with faintness or dizziness
    \item Heavy or painful bleeding after a long gap since the last period
    \item Sudden, severe abdominal or pelvic pain
    \item Chest pain, sudden breathlessness, or painful swelling of one leg, which may indicate a blood clot
    \item Signs of stroke: sudden facial drooping, arm weakness, or difficulty speaking
\end{itemize}

\section*{Diagnosis and Investigations}
In most cases the diagnosis of menopause is made clinically from the history of amenorrhoea and typical symptoms, and no tests are required for women over the age of 45. Investigations are used mainly to exclude other causes or to investigate abnormal bleeding.

\begin{itemize}
    \item \textbf{History:} menstrual pattern, symptom diary, timing of symptoms, and impact on sleep and daily life
    \item \textbf{Examination:} blood pressure, body mass index, and pelvic examination when bleeding symptoms or gynaecological concerns are present
    \item \textbf{Blood tests:} follicle-stimulating hormone (FSH) measurement may help confirm menopause in younger women or those in whom the diagnosis is unclear; it is not generally needed over age 45
    \item \textbf{Thyroid function tests and full blood count:} to exclude thyroid disease and anaemia as causes of symptoms
    \item \textbf{Pelvic ultrasound:} for heavy or irregular bleeding, to assess the endometrium, uterus, and ovaries
    \item \textbf{Endometrial sampling:} in post-menopausal bleeding, to exclude endometrial hyperplasia or cancer
\end{itemize}

\section*{Management}
Management is individualised and combines lifestyle measures with pharmacological treatment where symptoms warrant it. Treatment should be reviewed regularly as needs change.

\begin{itemize}
    \item \textbf{Lifestyle measures:} layered clothing, cooling techniques, regular exercise, and avoidance of alcohol, caffeine, and spicy food before sleep
    \item \textbf{Hormone replacement therapy (HRT):} the most effective treatment for vasomotor symptoms; oestrogen is used alone in women without a uterus and combined with progestogen in those with one, to protect the endometrium
    \item \textbf{Vaginal oestrogen:} low-dose topical therapy for genitourinary symptoms, with minimal systemic absorption
    \item \textbf{Non-hormonal medication:} some antidepressants and other agents can reduce hot flushes in women who cannot take or prefer to avoid HRT
    \item \textbf{Talking therapies:} cognitive behavioural therapy and mindfulness have evidence for improving vasomotor symptoms, sleep, and mood
    \item \textbf{Bone health:} adequate calcium and vitamin D, weight-bearing exercise, and bone density assessment where risk factors exist
    \item \textbf{Complementary approaches:} some women find benefit from herbal preparations, though evidence of effectiveness is limited and products are not tightly regulated
\end{itemize}

\section*{Complications}
\begin{itemize}
    \item Osteoporosis and increased fracture risk resulting from accelerated bone loss in the post-menopausal years
    \item Increased cardiovascular risk, including coronary heart disease and stroke, after menopause
    \item Genitourinary symptoms of menopause that affect intimacy, continence, and quality of life
    \item Chronic sleep disturbance and adverse effects on mood, work performance, and relationships
    \item Potential HRT side effects, including a small increased risk of venous thromboembolism, stroke, breast cancer, and gallbladder disease in some women, which must be weighed against benefits in shared decision-making
\end{itemize}

\section*{Prognosis and Outcomes}
Most women find that the most troublesome menopausal symptoms settle over time. Vasomotor symptoms typically improve within 2--5 years of the final period, although a minority of women continue to experience flushes for a decade or longer. HRT is highly effective at relieving vasomotor and many other symptoms while it is used. The long-term consequences of oestrogen deficiency, particularly bone loss and cardiovascular risk, persist after menopause and respond to lifestyle measures and, where appropriate, pharmacological bone and vascular protection. Women who experience premature or early menopause face greater cumulative health risks and are generally advised to consider hormone therapy until at least the natural age of menopause, with the balance of risks and benefits discussed individually.

\section*{Prevention and Self-Care}
\begin{itemize}
    \item Discuss symptoms openly with a doctor rather than accepting them as an unavoidable burden
    \item Maintain regular physical activity, including weight-bearing and resistance exercise
    \item Eat a balanced diet rich in calcium and vitamin D, and consider supplements if intake is inadequate
    \item Avoid smoking and limit alcohol, both of which can trigger flushes and worsen bone and cardiovascular health
    \item Use practical cooling measures: fans, cool drinks, layered clothing, and a cool bedroom
    \item Protect sleep with a consistent bedtime routine and treat any sleep disorder early
    \item Seek support for low mood, anxiety, or relationship strain, including talking therapies
    \item Consider bone density screening if there are risk factors for osteoporosis
\end{itemize}

\section*{Sources and Further Reading}
The following reputable sources provide further information for healthcare students and patients seeking reliable, current advice about the menopause and its management.

\begin{itemize}
    \item NHS. \textit{Menopause: symptoms, diagnosis, and treatment options.} https://www.nhs.uk/conditions/menopause/
    \item Mayo Clinic. \textit{Patient guide to menopause and hormone therapy.} https://www.mayoclinic.org/diseases-conditions/menopause/
    \item MSD Manual. \textit{Menopause: professional overview and management.} https://www.msdmanuals.com/
    \item MedlinePlus. \textit{Menopause health information.} https://medlineplus.gov/menopause.html
    \item World Health Organization. \textit{Menopause and women's health.} https://www.who.int/
    \item NICE. \textit{Menopause: diagnosis and management guidance.} https://www.nice.org.uk/guidance/ng23
\end{itemize}
